\scanpage{91}
% Typescript read after a half turn.
disjoints. \unsure{Donc|Dnnc} la donnée d'un homomorphisme de groupes $G\qquad\underline G_{m\ S}$ équivaut à la donnée d'une décomposition de $S$ comme réunion d'ouverts $S_m$ ($m\quad M$) deux à deux disjoints, i.e. à la donnée d'une application localement constante de $S$ dans $M$. Cela établit 1.3. donc 1.2.

\uline{Corollaire} 1.4. Un groupe diagonalisable est réflexif, donc aussi un groupe localement diagonalisable. Si $M,N$ sont deux groupes commutatifs ordinaires, \deleted{\unsure{on a}} l'homomorphisme naturel
\[
 \operatorname{Hom}_{S\text{-gr}}(M_S,N_S)\qquad\operatorname{Hom}_{S\text{-gr}}(D(N_S),D(M_S))
\]
est bijectif.

Notons maintenant que si $M$ est un groupe commutatif \uline{de type fini}, \deleted{\unsure{donc|fini}} donc défini par un nombre fini de générateurs et un nombre fini de relations, alors pour tout \deleted{préschéma en groupes $G$ sur $S$, on a une} \inserted{groupe ordinaire $N$,} l'application naturelle

On en déduit un isomorphisme de foncteurs sur $S$\unsure{:|=}
\[
 \underline{\operatorname{Hom}}_{S\text{-gr}}(M_S,N_S)\qquad\underline{\operatorname{Hom}}_{S\text{-gr}}(D(N_S),D(M_S))
\]
D'autre part \unsure{on|un} définit de façon évidente un homomorphisme
\[
 \operatorname{Hom}_{\mathrm{gr}}(M,N)_S\qquad\underline{\operatorname{Hom}}_{S\text{-gr}}(M_S,N_S)\ ,
\]
qui est toujours \deleted{\unsure{inj.}} un monomorphisme, et qui est un isomorphisme dès que $M$ est de type fini. On trouve ainsi:

\uline{Corollaire} 1.5. \deleted{\unsure{Supposons} $M$} Soient $M,N$ deux groupes commutatifs ordinaires, $M$ de type fini, alors \deleted{\unsure{on a}} $\underline{\operatorname{Hom}}_{S\text{-gr}}(D(N_S),D(M_S))$ est représentable, et on a un isomorphisme
\[
 \underline{\operatorname{Hom}}_{\mathrm{gr}}(M,N)_S\qquad\underline{\operatorname{Hom}}_{S\text{-gr}}(D(N_S),D(M_S))
\]
On en conclut plus généralement que si $G,H$ sont localement diagonalisables, $H$ étant de type fini, alors $\underline{\operatorname{Hom}}_{S\text{-gr}}(G,H)$ est représentable.
% archive sheet number: 90
